\par\addvspace{12pt}\Needspace{9\baselineskip}

\section{迭代产生的数列}

\label{sol:sec:analysis:8}

% ITEM: analysis-fixed-point-proof
\solutionitem{analysis-fixed-point-proof}{定理1.16的证明（1.8.2）}
设\(a\)是\(f \left( x \right)\)的不动点，函数\(f\)在\(a\)的邻域\(\left( a - r , a + r \right)\)中连续且严格单调递增，且在\(\left( a - r , a \right)\)上有\(f \left( x \right) > x\)，而区间\(\left( a , a + r \right)\)上有\(f \left( x \right) < x\)，那么由迭代产生的数列\(x_{n + 1} = f \left( x_{n} \right)\)只要第一项在\(\left( a - r , a + r \right)\)内，且不等于\(a\)，则以后就不会越出这个区间，而且是以\(a\)为极限的严格单调数列。
\solutionanswer
若\(a-r<x_1<a\)，由严格递增性及\(f(a)=a\)，有\(x_1<f(x_1)<a\)。归纳可知\(x_n<x_{n+1}<a\)，故各项都留在\([x_1,a]\)内，并严格递增。设极限为\(L\)，则\(x_1\le L\le a\)。由连续性在递推式中取极限，得\(L=f(L)\)。如果\(L<a\)，题设又给出\(f(L)>L\)，矛盾，故\(L=a\)。

若\(a<x_1<a+r\)，同理有\(a<x_{n+1}<x_n\)，数列严格递减且不越过\(a\)，极限仍只能为\(a\)。


% ITEM: analysis-iteration-exercises
\solutionitem{analysis-iteration-exercises}{例1.24}
完成以下计算或证明：

    （1）设\(a > 0 , x_{1} = \sqrt{a} , x_{n + 1} = \sqrt{a + x_{n}} , n \in \mathbb{N}^{*}\)，求\(\lim_{n \to \infty} x_{n}\).

    （2）设\(a > 0 , x_{1} = \sqrt{a} , x_{n + 1} = \sqrt{a x_{n}} , n \in \mathbb{N}^{*}\)，求\(\lim_{n \to \infty} x_{n}\).

    （3）设\(x_{1} = b , x_{n + 1} = \frac{1}{2} \left( x_{n}^{2} + 1 \right)\)，问：\(b\)取何值时数列\(\left\{ x_{n} \right\}\)收敛，并求其极限。

    （4）设\(\left\{ x_{n} \right\}\)为正数列，且满足\(x_{n + 1} + \frac{1}{x_{n}} < 2\)，证明：\(\left\{ x_{n} \right\}\)收敛，并求其极限。

    （5）设\(A > 0 , x_{1} > 0 , x_{n + 1} = \frac{x_{n} \left( x_{n}^{2} + 3 A \right)}{3 x_{n}^{2} + A} , n \in \mathbb{N}^{*}\)，证明：\(\left\{ x_{n} \right\}\)收敛于\(\sqrt{A}\)。
\solutionanswer
（1）不动点的正根为\(r=(1+\sqrt{1+4a})/2\)，满足\(r^2=a+r\)，故\(\sqrt a<r\)。函数\(\sqrt{a+x}\)严格递增，由\(x_1<x_2\)及归纳法得数列严格递增；又由\(x_1<r\)得以后各项都小于\(r\)。先用单调有界原理保证收敛，再取极限得到极限为\(r\)。

（2）作比例\(y_n=x_n/a\)，递推化为\(y_{n+1}=\sqrt{y_n}\)。若\(a>1\)，则\(y_1=1/\sqrt a<1\)，数列递增且小于\(1\)；若\(0<a<1\)，数列递减且大于\(1\)；\(a=1\)时为常数列。前两种情形的正极限满足\(L=\sqrt L\)，故\(L=1\)。因此\(x_n\to a\)。

（3）\(x_{n+1}-x_n=(x_n-1)^2/2\ge0\)。从第二项起各项均不小于\(1/2\)。若\(b\in[-1,1]\)，则\(x_2\in[1/2,1]\)，以后各项留在该区间并递增，极限满足\(L=(L^2+1)/2\)，故为\(1\)。若\(|b|>1\)，则\(x_2>1\)，尾列严格递增；若它有有限极限，也只能为\(1\)，与\(x_n\ge x_2>1\)矛盾，故趋于\(+\infty\)。所以收敛当且仅当\(b\in[-1,1]\)，极限为\(1\)。

（4）由下一项为正知\(2-1/x_n>0\)，故\(x_n>1/2\)。又
\[
x_{n+1}<2-\frac1{x_n}\le x_n,
\]
末步由\(x_n+1/x_n\ge2\)得到。因此数列严格递减且有正下界，设极限为\(L\ge1/2\)。取极限得\(L+1/L\le2\)，而均值不等式给出反向不等式，故只能\(L=1\)。

（5）按目标不动点作伸缩，令\(t_n=x_n/\sqrt A\)，则
\[
t_{n+1}=\frac{t_n(t_n^2+3)}{3t_n^2+1},\quad
 t_{n+1}-1=\frac{(t_n-1)^3}{3t_n^2+1},\quad
 t_{n+1}-t_n=\frac{2t_n(1-t_n^2)}{3t_n^2+1}.
\]
前一差式表明迭代不会越过\(1\)，后一差式表明两侧各自朝\(1\)移动。因此\(t_1<1\)时递增有上界，\(t_1>1\)时递减有下界，\(t_1=1\)时恒为\(1\)。极限为正，代入递推式得\(L=1\)，故\(x_n\to\sqrt A\)。
\par
